\documentclass{report}
\usepackage{amsmath,amssymb}
\setlength{\parindent}{0mm}
\setlength{\parskip}{1em}
\begin{document}
\begin{center}
\rule{6in}{1pt} \
{\large Luc Berger-Vergiat \\
{\bf Parallel domain decomposition preconditioners for Isogeometric discretizations arising in impact problems at high strain rates}}

Department of Civil Engineering \& Engineering Mechanics - room 626A \\ 610 Seeley W Mudd \\ 500 West 120th Street \\ New York \\ NY 10027
\\
{\tt lb2653@columbia.edu}\\
Luc Berger-Vergiat\end{center}

Shear bands are one of the most fascinating instabilities in metals that
occur under high strain rates. They describe narrow regions, on the order
of micron scales, where plastic deformations
and significant heating is localized that eventually leads to fracture
nucleation and failure of the component.

In this work shear bands are described by a set of four strongly coupled
thermo-mechanical equations discretized by a mixed finite element
formulation. A thermo-viscoplastic flow rule is used to model the
inelastic constitutive law and finite thermal conductivity is prescribed
which gives rise to an inherent physical length scale, governed by
competition of shear heating and thermal diffusion. The residual
equations are solved monolithically by a Newton type method at every time
step and has been shown to yield mesh insensitive result. The Jacobian of
the system is sparse and has a non-symmetric block structure that varies
with the different stages of shear bands formation.

The aim of this work is to develop fast converging Isogeometric elements
for monolithic solution of shear bands and appropriate parallel and
iterative solvers that are robust through all the deformation stages:
homogeneous, onset of instability and stress collapse. To this end, we
propose irreducible and hybrid NURBS quadrilateral elements that are
stable and locking free and lead to high rates of convergence. To solve
the resulting linearized systems, a novel Schur based
domain-decomposition preconditioners are proposed based on
constitutive/conservation laws splitting. The elements and solvers are
implemented in parallel and shows excellent scaling performance compared
to standard solvers.

Keywords: Shear bands, NURBS, Preconditioner, Schur complement, Domain
decomposition, Additive Schwarz method


\end{document}
